\begin{example}[Checking which colors separate a pair]
\label{ex:algorithm-separation-20260918}
Use the coloring drawn in \AlgoBlockRef{draft} from
\cref{ex:algorithm-submissions-20260918} and the forest in
\cref{ex:algorithm-peel-centroids-20260918}.  For the tested pair
$(i,j)=(1,10)$, the readout has neighbors
$\mathcal N_{\ResidualForest}(10)=\{1,11\}$.  The two tests in
\eqref{app:eq:pairwise-local-isolation} are therefore
\[
 \{1,11\}\cap\ColorClass{\tau}{2}=\emptyset,\qquad
 \{1,11\}\cap\ColorClass{\tau}{1}=\{1\}\subseteq\{i\}.
\]
Both hold: $h_\tau(1)=1$, $h_\tau(10)=2$, and $h_\tau(11)=3$.
If instead $h_\tau(11)=2$, then neighbor $11$ shares the readout color
and the first intersection contains $11$.  If instead $h_\tau(11)=1$,
then an additional neighbor shares the source color and the second
intersection is $\{1,11\}\not\subseteq\{1\}$.
This checks one pair and one coloring.  It does not assert the
simultaneous strict-majority event \eqref{app:eq:screen-success-event},
which requires the full family of $\ColoringCount$ colorings.
\end{example}
